\chapter{OPTIMASI BETON BERTULANG PADA STRUKTUR PORTAL RUANG}
\label{bab:optimasi}

\begin{quote}
\small\itshape
Catatan transkripsi: persamaan kendala \eqref{eq:3-1}--\eqref{eq:3-14} pada
bab ini juga aslinya berupa objek \textit{Equation Native} (MathType).
Bentuk persis setiap persamaan di bawah ini disusun ulang dengan mencocokkan
teks naratif (yang menyebut nama variabel setiap kendala secara eksplisit)
terhadap implementasi C++ yang mengimplementasikannya persis
(\texttt{balok::analisa()}, \texttt{balok::sengkang\_balok()},
\texttt{lendutan()} pada \texttt{Balok.hpp}; \texttt{kolom::rho()},
\texttt{kolom::analisa()}, \texttt{kolom::jarak\_tulangan()},
\texttt{kolom::kelangsingan()}, \texttt{kolom::sengkang\_kolom()} pada
\texttt{Kolom.hpp}), sehingga isi matematisnya terjamin sama dengan yang
benar-benar dihitung oleh program, walaupun nomor urut persamaan yang
persis seperti naskah asli tidak dapat dipastikan 100\%.
\end{quote}

\section{Struktur Program}

Program optimasi beton bertulang pada struktur portal ruang ini ditulis
dengan menggunakan bahasa pemrograman C++. Kode Program ini ditulis pada
beberapa file yaitu file utama yang berekstensi \texttt{cpp}, dan file
pembantu berekstensi \texttt{hpp}. Di dalam file yang berekstensi
\texttt{cpp} terdapat suatu fungsi bernama \texttt{main()}, fungsi ini akan
dijalankan pertama kali oleh program, selanjutnya dari dalam fungsi ini
akan dipanggil fungsi-fungsi lainnya sehingga membentuk suatu program yang
utuh.

Struktur program optimasi beton bertulang pada struktur portal ruang ini
merupakan gabungan dari pemrograman terstruktur dan pemrograman
berorientasi obyek.

Pemrograman terstruktur memiliki kelebihan dalam pengorganisasian dalam
membuat kode-kode program sehingga mudah dibaca dan mudah untuk dipahami,
gagasan pemrograman terstruktur pertama kali dikemukakan oleh Profesor
Edsger W. Dijkstra dari University of Eindhoven Nederland pada tahun 1965
(Sutedjo dan Michael, 1997, halaman 1). Dalam tulisannya, beliau menyatakan
bahwa pernyataan GOTO seharusnya jangan dipergunakan, namun pernyataan ini
ditanggapi oleh HD. Mills yang beranggapan bahwa pemrograman terstruktur
tidak hanya dihubungkan dengan pernyataan GOTO saja, tetapi pada struktur
program itu sendiri yang menentukan apakah program tersebut terstruktur
atau tidak, baik menggunakan pernyataan GOTO maupun tidak.

Pemrograman berorientasi obyek diciptakan karena masih dirasakan adanya
keterbatasan pada bahasa pemrograman tradisional, konsep pemrograman
berorientasi obyek yaitu membagi masalah-masalah ke dalam obyek, sehingga
data dan fungsi-fungsi yang akan dioperasikan digabung menjadi satu
kesatuan disebut pengkapsulan (\textit{encapsulated}).

Dalam bidang ilmu teknik, pemrograman terstruktur lebih banyak digunakan
karena rumitnya masalah yang dikerjakan dan memerlukan kemudahan
pemeriksaan (\textit{debugging}), dalam program ini struktur program utama
menggunakan konsep pemrograman terstruktur, sedangkan konsep pemrograman
berorientasi obyek digunakan pada analisis balok dan analisis kolom. Kelas
balok dan kelas kolom akan dibangkitkan apabila diperlukan analisa balok
atau analisa kolom dan apabila tidak diperlukan lagi kelas-kelas itu
dihancurkan, sehingga penggunaan kelas ini dapat menghemat memori komputer.

File yang terdapat pada program optimasi beton bertulang pada struktur
portal ruang adalah:
\begin{description}[style=nextline]
\item[\texttt{ORCISF.CPP}] merupakan file yang memuat fungsi \texttt{main()}.
\item[\texttt{BORLANDC.HPP}] file yang memuat header pustaka yang dibuat oleh Borland.
\item[\texttt{HEADER.HPP}] berisi kumpulan header.
\item[\texttt{PROTO.HPP}] berisi prototipe program.
\item[\texttt{VARIABEL.HPP}] berisi deklarasi variabel global.
\item[\texttt{INOUT.HPP}] berisi kumpulan sub program penanganan masalah masukan dan keluaran.
\item[\texttt{STRUKTUR.HPP}] berisi kumpulan sub program analisa struktur dengan metoda kekakuan.
\item[\texttt{SOLVER.HPP}] berisi sub program pembantu penanganan matrik.
\item[\texttt{KOLOM.HPP}] berisi kelas kolom untuk analisa kolom biaksial bertulangan simetri.
\item[\texttt{BALOK.HPP}] berisi kelas balok untuk analisa balok persegi bertulangan rangkap.
\item[\texttt{ELEMEN.HPP}] berisi kumpulan sub program yang menangani elemen pada balok dan kolom.
\item[\texttt{POLYHEDRON.HPP}] berisi sub program untuk optimasi dengan menggunakan metoda polihedron fleksibel.
\item[\texttt{PENORMALAN.HPP}] berisi kumpulan sub program yang menangani perubahan variabel balok dan kolom menjadi variabel normal.
\item[\texttt{PENGACAKAN.HPP}] berisi sub program yang mengacak nilai variabel.
\item[\texttt{DISKRITISASI.HPP}] berisi kumpulan sub program yang merubah nilai kontinu menjadi diskrit.
\item[\texttt{KENDALA.HPP}] berisi sub program untuk menghitung besarnya kendala yang terjadi pada struktur.
\item[\texttt{TELUSUR.HPP}] berisi sub program untuk melakukan penelusuran variabel.
\item[\texttt{BARU.HPP}] berisi sub program untuk menentukan struktur baru dalam metoda optimasi.
\item[\texttt{PENGURUTAN.HPP}] berisi sub program untuk mengurutkan fitness struktur.
\item[\texttt{TAMPILAN.HPP}] berisi sub program yang menangani masalah tampilan.
\item[\texttt{CETAK.HPP}] berisi sub program yang menangani masalah pencetakan hasil akhir.
\end{description}

\section{Sub Program Masukan dan Keluaran}

\subsection{Sub Program Untuk Menampilkan Menu Utama}

Pada saat program dijalankan pertama kali akan muncul sebuah menu yang
berisi pilihan untuk menjalankan program, menu utama program optimasi
beton bertulang pada struktur portal ruang seperti diperlihatkan pada
gambar 3-1.

\begin{center}
\fbox{\parbox{0.85\textwidth}{\footnotesize\ttfamily\raggedright\noindent
*************************************************************************\\
\hphantom{x}      Program Optimasi Beton Bertulang Pada Struktur Portal Ruang\\
\hphantom{xxxxxxxxxxxxxxxxxxxxxxxxx}Oleh : Yohan Naftali\\
\hphantom{xxxxxxxxxxxxxxxxxxxxxxxxxxxxxx}7712/TS\\
*************************************************************************\\
\hphantom{xxxxxxxxxxxxxxxxxxxxxxxxxxxxxxxxxxxxxxxxxxxxxxxxxxxxx}Time : 23:09:48.10\\
\hphantom{x} 1. Input data awal ke file\\
\hphantom{x} 2. Input data beban ke file\\
\hphantom{x} 3. Melihat isi file input\\
\hphantom{x} 4. Mengoptimasi Struktur\\
\hphantom{x} 5. Keluar\\
\hphantom{x} Pilihan (1-5) = \_
}}
\captionof{figure}{Menu Utama Program (tangkapan layar konsol asli, direproduksi sebagai teks)}
\end{center}

Pilihan yang terdapat pada menu utama terdiri dari:
\begin{enumerate}[label=\arabic*.]
\item Input data awal ke file, untuk memasukkan data tentang informasi struktur.
\item Input data beban ke file, untuk memasukkan data beban yang bekerja pada struktur.
\item Melihat isi file input, untuk menampilkan data yang telah dimasukkan.
\item Mengoptimasi struktur, untuk memulai proses optimasi beton bertulang pada struktur portal ruang.
\item Keluar, untuk menghentikan program.
\end{enumerate}
Fungsi-fungsi yang digunakan untuk menampilkan menu utama yang terdapat
pada file \texttt{TAMPILAN.HPP} terdiri dari:
\begin{enumerate}[label=\arabic*.]
\item \texttt{void menu\_utama()}, fungsi ini dipanggil dari fungsi
      \texttt{void main()}, selanjutnya dari fungsi ini dipanggil
      fungsi-fungsi lainnya.
\item \texttt{void about()}, fungsi ini dipanggil untuk menampilkan
      informasi program pada layar.
\end{enumerate}

\subsection{Sub Program Memasukkan Data}

Untuk melakukan proses pemasukkan data (\textit{input}), dapat dilakukan
dengan dua cara yaitu:
\begin{enumerate}[label=\arabic*.]
\item Mengikuti panduan yang disediakan oleh program.
\item Mengedit teks ASCII dengan bantuan MS-DOS Editor atau Notepad.
\end{enumerate}
Memasukkan data melalui panduan program, maka pertama kali akan diminta
nama file generik (tanpa ekstensi), nama file ini digunakan program untuk
memberi nama file-file bentukan lainnya. Selanjutnya program meminta
data-data struktur, mengedit teks ASCII dengan bantuan ASCII editor lebih
mudah digunakan untuk mengedit data masukan yang telah ada.

Sebelum memulai pemasukan data ke komputer, maka data-data struktur harus
kita persiapkan terlebih dahulu supaya mempermudah proses pemasukan data,
data-data struktur yang perlu kita persiapkan adalah:
\begin{enumerate}[label=\arabic*.]
\item Nama struktur
\item Jumlah batang pada struktur
\item Jumlah titik kumpul
\item Jumlah pengekang tumpuan pada struktur
\item Jumlah titik kumpul yang dikekang
\item Koordinat titik kumpul
\item Alokasi elemen batang pada struktur
\item Pengekang titik kumpul
\item Kuat desak beton karakteristik
\item Kuat tarik tulangan utama
\item Kuat tarik tulangan sengkang
\item Beban yang mengenai struktur
\item Data lebar dan tinggi balok
\item Data sisi penampang kolom
\item Data diameter tulangan utama
\item Data diameter tulangan sengkang
\item Data jumlah tulangan
\item Data jarak antara tulangan sengkang
\end{enumerate}
Setelah data kita masukkan ke dalam komputer, selanjutnya diperiksa lagi
data yang telah kita masukkan untuk memastikan bahwa data yang telah
dimasukkan sudah benar, apabila ada kesalahan, maka data masukan dapat
dikoreksi dengan mengedit file masukan.

Struktur file masukan pada program optimasi beton bertulang pada struktur
rangka ini adalah:
\begin{enumerate}[label=\arabic*.]
\item Data umum struktur terdapat pada file \texttt{generik.inp}
\item Data beban yang bekerja pada struktur terdapat pada file \texttt{generik.bbn}
\item Data diskrit mengenai ukuran penampang batang terdapat pada file \texttt{generik.isd}
\item Data diskrit mengenai diameter tulangan utama terdapat pada file \texttt{generik.idl}
\item Data diskrit mengenai diameter tulangan sengkang terdapat pada file \texttt{generik.ids}
\item Data diskrit mengenai jumlah tulangan utama pada salah satu penampang terdapat pada file \texttt{generik.ijl}
\item Data diskrit mengenai jarak antara tulangan sengkang terdapat pada file \texttt{generik.ijs}
\end{enumerate}
Fungsi-fungsi yang digunakan dalam proses masukan adalah:
\begin{enumerate}[label=\arabic*.]
\item \texttt{void input\_data()}, terdapat pada file \texttt{INOUT.HPP}, untuk menampilkan menu masukan
\item \texttt{void input\_data\_umum()}, terdapat pada file \texttt{INOUT.HPP}, untuk memasukkan data umum struktur ke file
\item \texttt{void input\_data\_diskrit()}, terdapat pada file \texttt{INOUT.HPP}, untuk memasukkan data elemen batang
\item \texttt{void load\_data()}, terdapat pada file \texttt{PEMBEBANAN.HPP}, untuk memasukkan data beban yang bekerja pada struktur
\end{enumerate}

\subsection{Sub Program Untuk Membaca Data Masukan}

Setelah data masukan diproses oleh komputer, dan diletakkan ke dalam
file-file masukan, maka untuk dapat dilakukan proses optimasi maka data
masukan tersebut perlu dibaca ulang, fungsi-fungsi yang digunakan untuk
melakukan proses ini adalah:
\begin{enumerate}[label=\arabic*.]
\item \texttt{void baca\_data()}, terdapat pada file \texttt{INOUT.HPP}, untuk membaca data masukan
\item \texttt{void baca\_beban()}, terdapat pada file \texttt{PEMBEBANAN.HPP}, untuk membaca data pembebanan.
\end{enumerate}

\section{Sub Program Inti}

\subsection{Sub Program Analisa Struktur}

Untuk menganalisa struktur portal ruang digunakan metoda kekakuan yang
dikembangkan oleh Weaver dan Gere (1980). Secara umum proses dalam
analisa struktur adalah sebagai berikut:
\begin{enumerate}[label=\arabic*.]
\item Pembentukan matrik rotasi
\item Pembentukan matrik kekakuan batang
\item Pembentukan vektor beban
\item Perhitungan perpindahan pada titik kumpul
\item Perhitungan gaya dan reaksi pada struktur
\end{enumerate}
Sub program analisa struktur terdapat pada file \texttt{STRUKTUR.HPP},
\texttt{PEMBEBANAN.HPP}, dan \texttt{SOLVER.HPP}. File \texttt{STRUKTUR.HPP}
berisi fungsi-fungsi untuk membentuk matrik rotasi, pembentukan matrik
kekakuan batang, perhitungan perpindahan, dan perhitungan gaya dan reaksi
pada struktur. Fungsi untuk membentuk vektor beban terdapat pada file
\texttt{PEMBEBANAN.HPP}, sedangkan file \texttt{SOLVER.HPP} terdapat
fungsi \texttt{void banfac()}, yang digunakan untuk melakukan faktorisasi
matrik simetris berjalur dengan pendekatan Choleski yang dimodifikasi, dan
terdapat pula fungsi \texttt{void bansol()}, yang digunakan untuk
mengolah matrik berjalur.

\subsection{Sub Program Analisa Balok Persegi Panjang}

Sub program analisa balok persegi panjang terdapat pada file
\texttt{BALOK.HPP}. Dalam file \texttt{BALOK.HPP} terdapat sebuah kelas
yang bernama \texttt{class balok}. Kelas ini dibangkitkan untuk menghitung
kendala yang terdapat pada setiap balok pada struktur.

Persamaan perhitungan kendala-kendala balok yang terdapat pada kelas
balok adalah:

\begin{enumerate}[label=\arabic*.]
\item \textbf{Kendala rasio tulangan maksimum}
\begin{equation}
\mathit{kendala\_rho\_b} = \max\!\left(\frac{\rho}{\rho_b} - 1,\; 0\right)
\label{eq:3-1}
\end{equation}
pada persamaan di atas $\rho$ (\texttt{RHO}) adalah rasio penulangan,
$\mathit{kendala\_rho\_b}$ adalah kendala rasio tulangan maksimum yang
dibatasi oleh $\rho_b$ (\texttt{RHB}) yaitu rasio penulangan pada keadaan
luluh dikalikan $0{,}75$ (persamaan \eqref{eq:2-21}--\eqref{eq:2-22}).

\item \textbf{Kendala rasio penulangan minimum}
\begin{equation}
\mathit{kendala\_rho\_m} = \max\!\left(\frac{\rho_{min}}{\rho} - 1,\; 0\right),
\qquad \rho_{min} = \frac{1{,}4}{f_y}
\label{eq:3-2}
\end{equation}
pada persamaan di atas $\mathit{kendala\_rho\_m}$ adalah kendala rasio
penulangan minimum, $\rho_{min}$ (\texttt{RMIN}) adalah rasio penulangan
minimum.

\item \textbf{Kendala momen lentur}
\begin{equation}
\mathit{kendala\_M} = \max\!\left(\frac{M_u}{\phi M_n} - 1,\; 0\right)
\label{eq:3-3}
\end{equation}
pada persamaan di atas $\mathit{kendala\_M}$ merupakan kendala momen
lentur, $M_u$ adalah momen yang bekerja pada struktur, $\phi M_n$
(\texttt{FMU}) adalah momen yang dapat ditahan oleh balok (persamaan
\eqref{eq:2-16}/\eqref{eq:2-20}).

\item \textbf{Kendala lendutan}
\begin{equation}
\mathit{kendala\_lendutan} = \max\!\left(\frac{\mathit{Lendutan}}{\mathit{Lendutan}_{ijin}} - 1,\; 0\right)
\label{eq:3-4}
\end{equation}
pada persamaan di atas $\mathit{kendala\_lendutan}$ adalah kendala
lendutan, $\mathit{Lendutan}$ adalah lendutan yang terjadi,
$\mathit{Lendutan}_{ijin}$ adalah lendutan ijin maksimum (dihitung
menurut SK SNI T-15-1991-03 pasal 3.2.5, fungsi \texttt{lendutan()} pada
\texttt{Balok.hpp}).

\item \textbf{Kendala tulangan geser}
\begin{equation}
\mathit{kendala\_sb} = \max\!\left(\frac{Jarak\_S}{S_{maks}} - 1,\; 0\right)
\label{eq:3-5}
\end{equation}
pada persamaan di atas $\mathit{kendala\_sb}$ adalah kendala tulangan
geser, $Jarak\_S$ adalah jarak sengkang, $S_{maks}$ (\texttt{SmakS})
adalah jarak antar sengkang maksimum.
\end{enumerate}

\subsection{Sub Program Analisa Kolom Biaksial}

Sub program analisa kolom biaksial terdapat pada file \texttt{KOLOM.HPP}.
Dalam file \texttt{KOLOM.HPP} terdapat sebuah kelas yang bernama
\texttt{class kolom}. Kelas ini dibangkitkan untuk menghitung kendala
yang terdapat pada setiap kolom pada struktur.

Persamaan perhitungan kendala-kendala kolom yang terdapat pada kelas
kolom adalah:

\begin{enumerate}[label=\arabic*.]
\item \textbf{Kendala rasio tulangan maksimum}
\begin{equation}
\mathit{kendala\_r\_mak} = \max\!\left(\frac{\rho}{0{,}08} - 1,\; 0\right)
\label{eq:3-6}
\end{equation}
pada persamaan di atas $\rho$ (\texttt{RHO}) adalah rasio penulangan,
$\mathit{kendala\_r\_mak}$ adalah kendala rasio tulangan maksimum yang
dibatasi $0{,}08$ ($8\,\%$).

\item \textbf{Kendala rasio penulangan minimum}
\begin{equation}
\mathit{kendala\_r\_min} = \max\!\left(\frac{0{,}01}{\rho} - 1,\; 0\right)
\label{eq:3-7}
\end{equation}
pada persamaan di atas $\mathit{kendala\_r\_min}$ adalah kendala rasio
penulangan minimum yang dibatasi sebesar $0{,}01$ ($1\,\%$).

\item \textbf{Kendala gaya aksial minimum}
\begin{equation}
\mathit{kendala\_po} = \max\!\left(\frac{P_n}{P_o} - 1,\; 0\right),
\qquad P_o = \phi\left(0{,}85 f_c' \,\mathit{sisi}^2 + AS_{tot} f_y\right)
\label{eq:3-8}
\end{equation}
pada persamaan di atas $\mathit{kendala\_po}$ merupakan kendala gaya
aksial minimum atau disebut juga kendala eksentrisitas minimum, $P_n$
merupakan gaya aksial nominal dan $P_o$ merupakan batas gaya aksial yang
dapat dipikul oleh kolom dikalikan faktor reduksi kekuatan untuk kolom
berpengikat sengkang yaitu $0{,}8$.

\item \textbf{Kendala gaya aksial}
\begin{equation}
\mathit{kendala\_pn} = \max\!\left(\frac{P_n}{P_{n,coba}} - 1,\; 0\right)
\label{eq:3-9}
\end{equation}
pada persamaan di atas $\mathit{kendala\_pn}$ adalah kendala gaya aksial
nominal, $P_n$ merupakan gaya aksial nominal yang terjadi pada struktur,
$P_{n,coba}$ adalah gaya aksial yang mampu didukung oleh kolom,
$P_{n,coba}$ dicari bersama-sama $M_{n,coba}$ berdasarkan eksentrisitas
yang terjadi pada kolom dengan metoda \textit{false position} untuk
menentukan letak garis netral (fungsi \texttt{kolom::hitung\_kolom()} pada
\texttt{Kolom.hpp}).

\item \textbf{Kendala momen ekivalen biaksial}\\
Untuk $M_{nx} > M_{ny}$ digunakan persamaan:
\begin{equation}
\mathit{kendala\_mn} = \max\!\left(\frac{M_{ox}}{M_{n,coba}} - 1,\; 0\right)
\label{eq:3-11a}
\end{equation}
untuk $M_{nx} \le M_{ny}$ digunakan persamaan:
\begin{equation}
\mathit{kendala\_mn} = \max\!\left(\frac{M_{oy}}{M_{n,coba}} - 1,\; 0\right)
\label{eq:3-11b}
\end{equation}
pada persamaan di atas $\mathit{kendala\_mn}$ merupakan kendala momen
ekivalen biaksial, $M_{ox}$ merupakan momen ekivalen biaksial pada sumbu
$X$, $M_{oy}$ merupakan momen ekivalen biaksial pada sumbu $Y$ (persamaan
\eqref{eq:2-23}--\eqref{eq:2-24}), $M_{n,coba}$ merupakan momen nominal
pada arah sumbu biaksial yang ditinjau.

\item \textbf{Kendala jarak tulangan minimum}
\begin{equation}
\mathit{kendala\_tul} = \max\!\left(\frac{jarak_{min}}{\mathit{jarak\_antar\_tulangan}} - 1,\; 0\right)
\label{eq:3-12}
\end{equation}
pada persamaan di atas $\mathit{kendala\_tul}$ adalah kendala jarak antara
tulangan minimum, $jarak_{min}$ adalah jarak minimum yang diijinkan
($\max(1{,}5\,dia,\,40\text{ mm})$ dibalik: yang diambil adalah nilai
terkecil dari $1{,}5\,dia$ dan $40$ mm), $\mathit{jarak\_antar\_tulangan}$
adalah jarak antara tulangan utama.

\item \textbf{Kendala kelangsingan kolom}
\begin{equation}
\mathit{kendala\_kelangsingan} = \max\!\left(\frac{\mathit{rasio\_kelangsingan}}{22} - 1,\; 0\right),
\qquad \mathit{rasio\_kelangsingan} = \frac{K L}{\mathit{sisi}\sqrt{1/12}}
\label{eq:3-13}
\end{equation}
pada persamaan di atas $\mathit{kendala\_kelangsingan}$ merupakan kendala
kelangsingan pada kolom, $\mathit{rasio\_kelangsingan}$ merupakan rasio
kelangsingan pada kolom (dengan $K=0{,}5$), yang dibatasi minimal $22$
sesuai SK SNI T-15-1991-03 ($KL/r \le 22$).
\end{enumerate}

\subsection{Sub Program Metoda Optimasi Polihedron Fleksibel}

Untuk melakukan proses optimasi digunakan metoda polihedron fleksibel,
prosedur umum yang digunakan dalam program optimasi beton bertulang pada
struktur portal ruang dengan metoda polihedron fleksibel adalah:
\begin{enumerate}[label=\arabic*.]
\item Penentuan variabel untuk struktur generasi pertama, yang dilakukan
      dengan pengacakan, dalam program optimasi beton bertulang pada
      struktur portal ruang ini penentuan variabel terdapat pada file
      \texttt{PENGACAKAN.HPP} dalam melakukan pengacakan, sebuah struktur
      dipilih supaya menggunakan elemen yang paling besar nilainya
      supaya dalam salah satu struktur tersebut ada yang tidak memiliki
      kendala, hal ini dimaksudkan supaya dalam proses optimasi tidak
      terjebak kepada optimal lokal yang berkendala.
\item Pembangkitan generasi pertama, sesudah seluruh struktur
      dibangkitkan, nilai fitness, kendala, harga, dan variabel desain
      disimpan dalam array.
\item Seluruh struktur pada generasi pertama diurutkan berdasarkan
      fitnessnya, pengurutan ini terdapat pada file
      \texttt{PENGURUTAN.HPP}, metoda pengurutan yang digunakan adalah
      metoda \textit{Bubble Sort}, metoda ini berasal dari sifat
      gelembung air yang semakin ringan akan menuju ke permukaan air.
\item Penentuan arah baru untuk melakukan penelusuran, terdapat pada
      file \texttt{TELUSUR.HPP}.
\item Menggantikan struktur terjelek dengan struktur baru hasil
      penelusuran yang memiliki fitness lebih baik, apabila dalam
      penelusuran tidak ditemukan struktur yang lebih baik dari pada
      struktur terjelek, maka dilakukan penyusutan menuju struktur
      terbaik.
\item Selanjutnya dilakukan lagi pengurutan fitness, proses ini
      dilakukan berulang-ulang sampai pada suatu batas konvergensi.
\end{enumerate}

Formulasi untuk fungsi penalti yang digunakan dalam program optimasi
beton bertulang pada struktur portal ruang adalah \textit{exterior
penalty function}, proses penelusuran dimulai dari daerah tidak layak,
yang kemudian secara bertahap menuju daerah layak. Formulasi fungsi
penalti untuk menyelesaikan masalah meminimumkan $f(x)$ yang memenuhi
kendala $g_j(x) \le 0$, $j=1,2,3,\ldots,n_g$ adalah:
\begin{equation}
\phi(x) = f(x) + r \sum_{j=1}^{n_g} \max\bigl(g_j(x),\,0\bigr)
\label{eq:3-14}
\end{equation}
dalam persamaan di atas $\phi$ merupakan fungsi penalti, $f(x)$ merupakan
fungsi sasaran yang berupa harga struktur yang formulasinya terdapat pada
persamaan \eqref{eq:1-12}, $r$ parameter penalti, $n_g$ jumlah kendala,
dan $g_j(x)$ persamaan kendala yang telah dihitung dalam kelas balok dan
kelas kolom (jumlah seluruh $\mathit{kendala\_*}$ dari persamaan
\eqref{eq:3-1}--\eqref{eq:3-13}, dijumlahkan menjadi \texttt{kendala}
pada konstruktor \texttt{balok()}/\texttt{kolom()}). Nilai fitness yang
dipakai untuk mengurutkan struktur (\texttt{fitstr}, lihat
\texttt{Kendala.hpp}/\texttt{Polyhedron.hpp}) dihitung sebagai
$\mathit{finalti}/(\mathit{harga} + \mathit{finalti}\times\mathit{kendala})$,
bentuk kebalikan dari fungsi penalti \eqref{eq:3-14} sehingga struktur
termurah dan paling memenuhi kendala memberikan fitness tertinggi.

Nilai $r$ harus cukup besar agar mampu mengarahkan proses penelusuran
dekat dengan daerah layak, tetapi cukup kecil agar proses peminimuman
tidak mengalami kesulitan.

Dalam program optimasi beton bertulang pada struktur portal ruang
konvergensi ditentukan sebagai berikut:
\begin{enumerate}[label=\arabic*.]
\item Jumlah iterasi melebihi batas yang telah ditentukan
\item Terjadi penyusutan berturut yang menghasilkan fitness yang sama
      sebanyak jumlah struktur desain
\item Terdapat struktur berjumlah sama dengan jumlah variabel desain
      yang memiliki nilai fitness sama.
\end{enumerate}

\subsection{Sub Program Mencetak Hasil Akhir Optimasi}

Setelah suatu proses optimasi selesai, hasil optimasi berupa nilai-nilai
variabel desain terbaik, harga, kendala dan hasil perhitungan struktur
dicetak, sub program untuk mencetak hasil akhir tersebut ke dalam file
adalah \texttt{void cetak\_akhir()} yang terdapat pada file
\texttt{CETAK.HPP}.

File-file hasil keluaran yang dicetak oleh program optimasi beton
bertulang pada struktur portal ruang adalah:
\begin{enumerate}[label=\arabic*.]
\item \texttt{generik.opt} berisi dimensi balok dan kolom pada struktur
      yang paling optimal.
\item \texttt{generik.str} berisi hasil perhitungan gaya batang dengan
      metoda kekakuan.
\item \texttt{generik.kdl} berisi kendala yang terdapat pada struktur.
\item \texttt{generik.inf} berisi informasi data masukan.
\item \texttt{generik.his} berisi riwayat optimasi.
\end{enumerate}
